\abstract{Retrieval-Augmented Generation (RAG) for Arabic faces a memory bottleneck: Arabic's rich morphology yields more tokens than English text, inflating the Key-Value (KV) cache engines like vLLM must maintain per request. We introduce a Lightweight Semantic Pruning Middleware (LSPM) that scores retrieved passages at the sentence level with a cross-encoder, keeping only the most relevant subset at a fixed or dynamic ratio. We report five preliminary findings. First, across 30 verified Arabic/English sentence pairs, two tokenizers confirm a token ratio of 1.76x/1.69x (bootstrap p $<$ 0.0001 vs. parity). Second, a fidelity pilot (r=0.3-0.7) shows pruned answers match unpruned ones (ROUGE-L 0.89-0.93, BERTScore-F1 0.96-0.97). Third, on 140 ARCD questions (980 generations), LSPM beats naive truncation on F1 at the most aggressive ratio (+0.114, adjusted p = 0.005) and is equivalent to the unpruned answer (p = 0.038, 0.003); naive truncation is significantly worse than unpruned at two ratios (adjusted p = 0.0004, 0.015). Fourth, a GPU benchmark shows LSPM lowers KV-cache occupancy at every ratio at c = 1, with 8.5-11.9\% lower throughput and near-unchanged token rate, exploratory only since it reverses at higher concurrency. Fifth, two blind evaluators rated 90 answers for correctness and faithfulness (kappa 0.79, 0.94): LSPM scored at least as high as raw and higher than naive truncation, consistent with the F1 result above. We release the implementation, raw data, and a live demo.}

\keywords{retrieval-augmented generation, Arabic natural language processing, prompt compression, vLLM, large language models}
